\documentclass[11pt,a4paper]{article}
\usepackage[utf8]{inputenc}
\usepackage[T1]{fontenc}
\usepackage[french]{babel}
\usepackage{geometry}
\usepackage{xcolor}
\usepackage{hyperref}
\usepackage{titlesec}

\geometry{margin=1in}
\definecolor{titleblue}{RGB}{0,51,102}

\hypersetup{
    colorlinks=true,
    linkcolor=titleblue,
    urlcolor=titleblue,
    pdftitle={Lettre de Motivation - Ismail AISSAOUI},
}

\begin{document}

\noindent \textbf{Ismail AISSAOUI} \\
Ingénieur d'État en Intelligence Artificielle (ESI-SBA) \\
Email: ismail.aissaoui.pro@gmail.com \\
Téléphone: +213 660 70 77 96 \\
LinkedIn: https://linkedin.com/in/aissaoui-ismail-6a77a92a7
Portfolio: https://ismail-aissaoui.vercel.app

\bigskip

\begin{flushright}
    Au Comité de Sélection du Doctorat \\
    \textbf{Programme EDT - Projet Catalyst} \\
    À l'attention de : Prof. Fabrice Gamboa, Dr. Mathieu Vallee, Dr. Clément Gauchy \\
    \medskip
    1er mai 2026
\end{flushright}

\bigskip

\noindent \textbf{Objet : Candidature au doctorat : Apprentissage automatique interprétable et robuste pour les simulations numériques basées sur la physique (Rang 1 - PC1)}

\bigskip

Madame, Monsieur,

Ingénieur d'État en Intelligence Artificielle diplômé de l'ESI-SBA, je vous adresse avec enthousiasme ma candidature pour le poste de doctorat intitulé « Apprentissage automatique interprétable et robuste pour les simulations numériques basées sur la physique » au sein du programme national Engineering Digital Twin (EDT). Cette candidature est motivée par la forte adéquation de mon profil avec les objectifs scientifiques du projet Catalyst.

Actuellement en fin de cycle ingénieur, je réalise ma thèse de fin d'études chez Sonatrach, où j'ai développé un **Jumeau Numérique informé par la physique** pour la surveillance de flottes de turbines à gaz. Mes recherches se concentrent sur l'utilisation de l'apprentissage automatique pour émuler et accélérer les solveurs numériques complexes—un parallèle direct avec les objectifs du projet Catalyst (Reliable Hybrid Model Forge). Plus précisément, j'ai conçu des fonctions de perte personnalisées de type **PINN (Physics-Informed Neural Network)** pour garantir la cohérence thermodynamique au sein de modèles de substitution génératifs basés sur les architectures **Mamba-SSM** et **xLSTM**.

Le défi central de votre projet—traiter la nature « boîte noire » des modèles de substitution ML et leur sensibilité aux données hors distribution—est précisément ce qui motive mes travaux actuels. Dans le cadre de mon projet chez Sonatrach, j'ai été confronté au besoin critique d'**indicateurs quantitatifs de fiabilité** lors du déploiement de modèles de substitution sur des dispositifs industriels. Je suis particulièrement impatient d'appliquer des méthodes statistiques avancées, telles que l'**Analyse de Sensibilité**, le **Transport Optimal** et la **Profondeur Statistique**, pour définir les domaines de fiabilité de ces modèles hybrides.

Ma maîtrise de l'implémentation sur PyTorch et C++ en environnement de production, combinée à ma rigueur en statistiques et méthodes numériques, me permettront de contribuer activement à la **plateforme Artemis** et à l'écosystème EDT. Je suis un chercheur autonome, capable de faire le pont entre la théorie statistique et les simulations physiques industrielles.

Je vous remercie de l'attention que vous porterez à ma candidature. Je reste à votre entière disposition pour un entretien afin de discuter de la manière dont mon expérience en IA informée par la physique peut contribuer au succès du projet Catalyst.

Dans l'attente de votre réponse, je vous prie d'agréer, Madame, Monsieur, l'expression de mes salutations distinguées.

\bigskip

\noindent Ismail AISSAOUI

\end{document}
